% GENERATED FILE. Edit canonical lessons, then run npm run generate:books.
% canonical-source-hash: fnv1a64:18ad686361b61b48
% canonical-lessons: BN-C503-oyebsait, BN-C503-lingk, BN-C503-onlain, BN-C503-skrin, BN-C503-daunlod-kora

\chapter{Online}
\label{ch:bn-online}
\hlchaptermodality{503}

\begin{chapteropening}
\textbf{By the end of this chapter:} \emph{I can talk about websites, links and downloading.}

Complete the last lesson of chapter 503: I can talk about websites, links and downloading.
\end{chapteropening}

\section[\texorpdfstring{\bn{ওয়েবসাইট} (oyebsāiṭ)}{ওয়েবসাইট (oyebsāiṭ)}]{\bn{ওয়েবসাইট} (oyebsāiṭ) — a website}

\label{lesson:BN-C503-oyebsait}



\begin{quote}
Before the new one: say the Bengali for to rummage, then the Bengali for to charge a phone.
\end{quote}

\subsection*{You'll want to know: \bn{ওয়েবসাইট}}
\textbf{\bn{ওয়েবসাইট}} — \emph{oyebsāiṭ} — \textquotedblleft{}a website\textquotedblright{}.

\begin{quote}
\bn{স্কুলের} \bn{ওয়েবসাইটে} \bn{নোটিশটা} \bn{আছে।}
\end{quote}

\begin{quote}
\emph{Skuler oyebsāiṭe noṭishṭā āchhe.}
\end{quote}

\begin{quote}
The notice is on the school's website.
\end{quote}

\begin{cousinweb}
\textbf{\bn{ওয়েবসাইট}} (\emph{oyebsāiṭ}) is borrowed from English \emph{website}.
\end{cousinweb}

\begin{grammarlens}[title={What you've built}]
One more word for describing what you lived and felt.
\end{grammarlens}

\subsection*{Your turn}
\begin{itemize}
\raggedright
  \item \emph{Say it:} \emph{oyebsāiṭ}
  \item \emph{Say it:} \emph{oyebsāiṭ}, once more
  \item \emph{Say it:} \emph{Skuler oyebsāiṭe noṭishṭā āchhe.} — The notice is on the school's website.
  \item \emph{Recall:} say the Bengali for to rummage, then the Bengali for to charge a phone, then say \emph{oyebsāiṭ} again
\end{itemize}

\begin{tcolorbox}[breakable,colback=teal!4,colframe=teal!35!black,title={Before you move on}]
What does \bn{ওয়েবসাইট} mean? (\textquotedblleft{}A website\textquotedblright{}.) Say it once more.
\end{tcolorbox}
\section[\texorpdfstring{\bn{লিংক} (lingk)}{লিংক (lingk)}]{\bn{লিংক} (lingk) — a link online; a web link}

\label{lesson:BN-C503-lingk}



\begin{quote}
Before the new one: say the Bengali for to charge a phone, then the Bengali for a website.
\end{quote}

\subsection*{You'll want to know: \bn{লিংক}}
\textbf{\bn{লিংক}} — \emph{lingk} — \textquotedblleft{}a link online; a web link\textquotedblright{}.

\begin{quote}
\bn{আমি} \bn{বন্ধুকে} \bn{ভিডিওর} \bn{লিংক} \bn{পাঠালাম।}
\end{quote}

\begin{quote}
\emph{Āmi bôndhuke bhiḍior lingk pāṭhālām.}
\end{quote}

\begin{quote}
I sent my friend the link to the video.
\end{quote}

\begin{cousinweb}
\textbf{\bn{লিংক}} (\emph{lingk}) is borrowed from English \emph{link}.
\end{cousinweb}

\begin{grammarlens}[title={What you've built}]
One more word for describing what you lived and felt.
\end{grammarlens}

\subsection*{Your turn}
\begin{itemize}
\raggedright
  \item \emph{Say it:} \emph{lingk}
  \item \emph{Say it:} \emph{lingk}, once more
  \item \emph{Say it:} \emph{Āmi bôndhuke bhiḍior lingk pāṭhālām.} — I sent my friend the link to the video.
  \item \emph{Recall:} say the Bengali for to charge a phone, then the Bengali for a website, then say \emph{lingk} again
\end{itemize}

\begin{tcolorbox}[breakable,colback=teal!4,colframe=teal!35!black,title={Before you move on}]
What does \bn{লিংক} mean? (\textquotedblleft{}A link online; a web link\textquotedblright{}.) Say it once more.
\end{tcolorbox}
\section[\texorpdfstring{\bn{অনলাইন} (ônlāin)}{অনলাইন (ônlāin)}]{\bn{অনলাইন} (ônlāin) — online}

\label{lesson:BN-C503-onlain}



\begin{quote}
Before the new one: say the Bengali for a website, then the Bengali for a link online.
\end{quote}

\subsection*{You'll want to know: \bn{অনলাইন}}
\textbf{\bn{অনলাইন}} — \emph{ônlāin} — \textquotedblleft{}online\textquotedblright{}.

\begin{quote}
\bn{আজকাল} \bn{আমরা} \bn{অনলাইনে} \bn{টিকিট} \bn{কিনি।}
\end{quote}

\begin{quote}
\emph{Ājkāl āmrā ônlāine ṭikiṭ kini.}
\end{quote}

\begin{quote}
Nowadays we buy tickets online.
\end{quote}

\begin{cousinweb}
\textbf{\bn{অনলাইন}} (\emph{ônlāin}) is borrowed from English \emph{online}; it takes the locative: \bn{অনলাইনে}, online.
\end{cousinweb}

\begin{grammarlens}[title={What you've built}]
One more word for describing what you lived and felt.
\end{grammarlens}

\subsection*{Your turn}
\begin{itemize}
\raggedright
  \item \emph{Say it:} \emph{ônlāin}
  \item \emph{Say it:} \emph{ônlāin}, once more
  \item \emph{Say it:} \emph{Ājkāl āmrā ônlāine ṭikiṭ kini.} — Nowadays we buy tickets online.
  \item \emph{Recall:} say the Bengali for a website, then the Bengali for a link online, then say \emph{ônlāin} again
\end{itemize}

\begin{tcolorbox}[breakable,colback=teal!4,colframe=teal!35!black,title={Before you move on}]
What does \bn{অনলাইন} mean? (\textquotedblleft{}Online\textquotedblright{}.) Say it once more.
\end{tcolorbox}
\section[\texorpdfstring{\bn{স্ক্রিন} (skrin)}{স্ক্রিন (skrin)}]{\bn{স্ক্রিন} (skrin) — a screen}

\label{lesson:BN-C503-skrin}



\begin{quote}
Before the new one: say the Bengali for a link online, then the Bengali for online.
\end{quote}

\subsection*{You'll want to know: \bn{স্ক্রিন}}
\textbf{\bn{স্ক্রিন}} — \emph{skrin} — \textquotedblleft{}a screen\textquotedblright{}.

\begin{quote}
\bn{ল্যাপটপের} \bn{স্ক্রিনটা} \bn{বড়।}
\end{quote}

\begin{quote}
\emph{Læpṭôper skrinṭā bôṛo.}
\end{quote}

\begin{quote}
The laptop's screen is big.
\end{quote}

\begin{cousinweb}
\textbf{\bn{স্ক্রিন}} (\emph{skrin}) is borrowed from English \emph{screen}.
\end{cousinweb}

\begin{grammarlens}[title={What you've built}]
One more word for describing what you lived and felt.
\end{grammarlens}

\subsection*{Your turn}
\begin{itemize}
\raggedright
  \item \emph{Say it:} \emph{skrin}
  \item \emph{Say it:} \emph{skrin}, once more
  \item \emph{Say it:} \emph{Læpṭôper skrinṭā bôṛo.} — The laptop's screen is big.
  \item \emph{Recall:} say the Bengali for a link online, then the Bengali for online, then say \emph{skrin} again
\end{itemize}

\begin{tcolorbox}[breakable,colback=teal!4,colframe=teal!35!black,title={Before you move on}]
What does \bn{স্ক্রিন} mean? (\textquotedblleft{}A screen\textquotedblright{}.) Say it once more.
\end{tcolorbox}
\section[\texorpdfstring{\bn{ডাউনলোড} \bn{করা} (ḍāunloḍ kôrā)}{ডাউনলোড করা (ḍāunloḍ kôrā)}]{\bn{ডাউনলোড} \bn{করা} (ḍāunloḍ kôrā) — to download}

\label{lesson:BN-C503-daunlod-kora}



\begin{quote}
Before the new one: say the Bengali for online, then the Bengali for a screen.
\end{quote}

\subsection*{You'll want to know: \bn{ডাউনলোড} \bn{করা}}
\textbf{\bn{ডাউনলোড} \bn{করা}} — \emph{ḍāunloḍ kôrā} — \textquotedblleft{}to download\textquotedblright{}.

\begin{quote}
\bn{আমি} \bn{ইন্টারনেট} \bn{থেকে} \bn{একটা} \bn{গান} \bn{ডাউনলোড} \bn{করলাম।}
\end{quote}

\begin{quote}
\emph{Āmi inṭārneṭ theke ekṭā gān ḍāunloḍ kôrlām.}
\end{quote}

\begin{quote}
I downloaded a song from the internet.
\end{quote}

Past (he/she …): \bn{তিনি} \bn{ছবিটা} \bn{ডাউনলোড} \bn{করলেন}, \emph{tini chhôbiṭā ḍāunloḍ kôrlen}.

\begin{cousinweb}
\textbf{\bn{ডাউনলোড} \bn{করা}} (\emph{ḍāunloḍ kôrā}) joins \emph{download}, borrowed from English, to \bn{করা}, \textquotedblleft{}to do\textquotedblright{}.
\end{cousinweb}

\begin{grammarlens}[title={What you've built}]
One more word for describing what you lived and felt.
\end{grammarlens}

\subsection*{Your turn}
\begin{itemize}
\raggedright
  \item \emph{Say it:} \emph{ḍāunloḍ kôrā}
  \item \emph{Say it:} \emph{ḍāunloḍ kôrā}, once more
  \item \emph{Say it:} \emph{Āmi inṭārneṭ theke ekṭā gān ḍāunloḍ kôrlām.} — I downloaded a song from the internet.
  \item \emph{Recall:} say the Bengali for online, then the Bengali for a screen, then say \emph{ḍāunloḍ kôrā} again
\end{itemize}

\begin{tcolorbox}[breakable,colback=teal!4,colframe=teal!35!black,title={Before you move on}]
What does \bn{ডাউনলোড} \bn{করা} mean? (\textquotedblleft{}To download\textquotedblright{}.) Say it once more.
\end{tcolorbox}
